\section{Discussion}
\label{sec:discussion}

\subsection{Does This Close the Gap?}

Section~\ref{sec:background} lined up related work against three
requirements, automated, real-time, and on physical hardware, and
found that no prior approach met all three at once
(Table~\ref{tab:related_work}). This system computes trial time, path
length, and rate of orientation change directly from PSM kinematics,
checks cylinder orientation and object state directly from vision,
and reports contact force, all without a human rater, on the physical
dVRK, during the trial itself (Section~\ref{sec:methods}). The one
exception is 3D localization, which runs post-hoc, not live; but it
was never delivered as one of the real-time measures
(Section~\ref{subsec:localization}), it exists to validate the
others, so it is not itself a gap in that claim.
% TODO: revisit this once Results is complete, in case the timing
% numbers change this conclusion.

\subsection{What the Timing Means}
\label{subsec:discussion_timing}

% PLACEHOLDER, pending Section~\ref{subsec:results_timing}'s numbers.
% TODO: interpret the measured latencies. Is the slowest module fast
% enough that a trainee would not perceive the feedback as delayed,
% e.g. compared against typical human reaction time, or against how
% Laca et al.'s feedback loop was real-time even though the
% measurement behind it was not (Section~\ref{sec:background})? Which
% module is the bottleneck, and does closing it matter for the
% trainee-facing goal (Section~\ref{subsec:goal})?

\subsection{Limitations}

This thesis reports bench trials run by the author, not a trainee
study (Section~\ref{sec:results}); how these measures perform for
actual novice and expert trainees, the population this system is
ultimately for (Section~\ref{subsec:motivation}), is untested.

The object-state classifier's reliability is not yet established: the
rule-based labeler it was bootstrapped from struggled specifically to
tell the held and dropped states apart
(Section~\ref{subsec:object_state}), and the learned classifier's own
cross-validated accuracy is not yet measured. Because placement
scoring and object state are reported independently
(Section~\ref{subsec:object_state}), a classifier error here does not
corrupt the orientation-based measures, but it does limit how much the
object-state signal itself can be trusted for now.

3D localization only runs post-hoc, at roughly 5 frames per second
(Section~\ref{subsec:localization}), so it cannot itself be a
delivered real-time measure; an earlier, real-time version was
abandoned after it proved too laggy and inaccurate to trust. This is
a deliberate scope choice, not an oversight, but it does mean cylinder
position is not independently validated live, only after the fact.

Both the object-state classifier and 3D localization track one
cylinder at a time (Sections~\ref{subsec:object_state}
and~\ref{subsec:localization}), even though the task itself allows
multiple cylinders in play at once (Section~\ref{subsec:task});
extending them to the multi-cylinder case is not yet done.

The HSV thresholds behind cylinder orientation
(Section~\ref{subsec:angular_displacement}) were tuned for one
camera, one board, and one lighting setup; whether they still work
under different lighting or camera placement has not been tested.

\subsection{What This Sets Up}

Two of the formative study's findings
(Section~\ref{subsec:formative_study}) point past this thesis rather
than into it: a redesigned, lower-peg task board, built but not used
here, and the idea of a closed-loop, difficulty-adaptive task. Both,
along with a trainee study of the system itself, are future work,
taken up in Section~\ref{sec:conclusion}.
